\documentclass[11pt]{article}
\usepackage[utf8]{inputenc}
\usepackage{amsmath,amssymb}
\usepackage{geometry}
\geometry{margin=1in}

\title{S-wind: Autonomous Algebraic Discovery with Formal Discovery Criterion}
\author{Valhalla Research Team}
\date{January 2026}

\begin{document}
\maketitle

\section*{Summary}

We introduce \textbf{S-wind}, an autonomous algebraic discovery framework with two key innovations:

\begin{enumerate}
    \item \textbf{Extended Precision Rational (EPR)}: Arbitrary precision exact arithmetic with zero floating-point error
    \item \textbf{Formal Discovery Criterion}: A result is a \emph{discovery} iff it differs from known authoritative formulas AND is verified correct
\end{enumerate}

\section*{Key Results}

\begin{itemize}
    \item \textbf{17 experiments} completed in 2.4 seconds
    \item \textbf{8 known sequences} correctly recognized (Fibonacci, Lucas, Pell, etc.)
    \item \textbf{6 pattern detections} correctly classified as non-discoveries
    \item \textbf{3 impossibility theorems} automatically applied (primes, Catalan, partition)
    \item \textbf{2 novel discoveries}: Harmonic numerators (degree-7 polynomial), Basel partial sum (degree-4 polynomial)
\end{itemize}

\section*{Significance}

The formal discovery criterion distinguishes S-wind from existing CAS systems, which cannot distinguish between pattern fitting and genuine novelty.

\section*{Availability}

Source code: Valhalla/s-wind-sim (Rust, MIT License)

\end{document}
